%&17pt
\input инструкция

\sec Запуск оповещения

\subsec 1. Создание звукозаписи

\item{1)} Запустить программу записи.

$$\epsfbox{Запуск_оповещения/01-создать_звукозапись/01.eps}$$

\item{2)} Нажать микрофон, записать сообщение. В конце сообщения сказать
«ДЛЯ ПОДТВЕРЖДЕНИЯ НАЖМИТЕ ОДИН».

$$\epsfbox{Запуск_оповещения/01-создать_звукозапись/02.eps}$$

\item{3)} Остановить запись сообщения.

$$\epsfbox{Запуск_оповещения/01-создать_звукозапись/03.eps}$$

\item{4)} Прослушать последнюю запись (самая верхняя строка).
Если запись плохая, то записать новую.

$$\epsfbox{Запуск_оповещения/01-создать_звукозапись/04.eps}$$

\item{5)} Открыть папку «Звукозапись».

$$\epsfbox{Запуск_оповещения/01-создать_звукозапись/05.eps}$$

\item{6)} Щёлкнуть два раза на ogg с именем последней записи.

$$\epsfbox{Запуск_оповещения/01-создать_звукозапись/06.eps}$$

\item{7)} Вместо ogg должен появиться файл wav с таким же именем.

$$\epsfbox{Запуск_оповещения/01-создать_звукозапись/07.eps}$$

\subsec 2. Запуск оповещения

\item{1)} Запустить браузер либо выбрать открытый.

$$\epsfbox{Запуск_оповещения/02-запустить_оповещение/01.eps}$$

\item{2)} Запуск АСО (пропустить этот шаг если в первом пункте выбрали открытый браузер).

$$\epsfbox{Запуск_оповещения/02-запустить_оповещение/02.eps}$$

\item{3)} Вход в АСО (пропустить этот шаг если в первом пункте выбрали открытый браузер).

$$\epsfbox{Запуск_оповещения/02-запустить_оповещение/03.eps}$$

\item{4)} Выбрать в столбце слева «Группы оповещений».

$$\epsfbox{Запуск_оповещения/02-запустить_оповещение/04.eps}$$

\vfill\eject

\item{5)} Выбираем группу (щёлкаем по названию а не по квадратику слева от названия).

$$\epsfxsize=188mm \epsfbox{Запуск_оповещения/02-запустить_оповещение/05.eps}$$

\item{6)} Убедиться что была выбрана необходимая группа.

$$\epsfbox{Запуск_оповещения/02-запустить_оповещение/06.eps}$$

\item{7)} Загрузить файл звукозаписи, созданный как написано выше.

$$\epsfbox{Запуск_оповещения/02-запустить_оповещение/07.eps}$$

\item{8)} Выбрать файл.

$$\epsfbox{Запуск_оповещения/02-запустить_оповещение/08.eps}$$

\vfill\eject

\item{9)} Подтвердить выбор.

$$\epsfbox{Запуск_оповещения/02-запустить_оповещение/09.eps}$$

\item{10)} Сохранить.

$$\epsfbox{Запуск_оповещения/02-запустить_оповещение/10.eps}$$

\item{11)} Запустить оповещение.

$$\epsfbox{Запуск_оповещения/02-запустить_оповещение/11.eps}$$

\item{12)} Подтвердить запуск.

$$\epsfbox{Запуск_оповещения/02-запустить_оповещение/12.eps}$$

\vfill\eject

\subsec 3. Печать отчёта

\item{1)} Выбрать в столбце слева «Сеансы оповещений».

$$\epsfbox{Запуск_оповещения/03-печать_отчета/01.eps}$$

\item{2)} Сформировать отчёт.

$$\epsfbox{Запуск_оповещения/03-печать_отчета/02.eps}$$

\item{3)} Скачать PDF.

$$\epsfbox{Запуск_оповещения/03-печать_отчета/03.eps}$$

\item{4)} Открыть PDF.

$$\epsfbox{Запуск_оповещения/03-печать_отчета/04.eps}$$

\vfill\eject

\item{5)} Щёлкнуть «Файл» и в появившемся меню выбрать «Печать».

$$\epsfbox{Запуск_оповещения/03-печать_отчета/05.eps}$$

\item{6)} Распечатать отчёт.

$$\epsfbox{Запуск_оповещения/03-печать_отчета/06.eps}$$

\subsec 4. Очистка звукозаписи

\item{1)} Закрыть программу звукозаписи.

$$\epsfbox{Запуск_оповещения/04-очистить_звукозапись/01.eps}$$

\vfill\eject

\item{2)} Открыть папку «Звукозапись».

$$\epsfbox{Запуск_оповещения/04-очистить_звукозапись/02.eps}$$

\item{3)} Щёлкнуть правой кнопкой на каждом файле и в появившемся меню выбрать «Удалить».

$$\epsfbox{Запуск_оповещения/04-очистить_звукозапись/03.eps}$$

\bye
